\chapter{Related Work}
\label{ch:related}

The framework joins methods developed for different purposes. Active perception explains why motion can improve a measurement. Task-parameterized learning explains how a demonstration can be expressed in more than one frame. Convex programming provides a way to choose forces within explicit friction constraints. Compliant assembly explains why contact geometry and mechanical stiffness must be considered together. Their combination is useful only if the assumptions at the interfaces are preserved.

\section{Handling uncertain inertial properties}

The preprint by Wong et al. studies industrial dual-arm box handling with online inertial estimation and convex wrench optimization~\cite{wong2026}. Its motivation is directly relevant: packages can have a simple visible boundary and uncertain contents. The work estimates mass distribution from contact wrenches and formulates friction feasibility as a constraint. It also considers torsional contact capacity through a friction limit surface.

The present contact model differs at the level of representation. A finite patch can be assigned a resultant force and torsional moment with a coupled limit surface. Here, each hand has exactly three point forces. The wrist moment is generated by summing their moment arms. This avoids assigning an independent torsional moment to an individual point. It also makes loss of one pad explicit. The two descriptions may both represent useful contact physics, but their variables and admissible wrench sets are not interchangeable.

The connection with industrial handling is consequently specific. Online estimation reduces dependence on a pre-existing load model, and constrained allocation prevents the optimizer from requesting a force outside its stated contact set. Neither method removes the need to measure contact geometry or to verify that the commanded forces can be transmitted by the wrist and tool.

\section{Active perception and informative motion}

Bajcsy's active-perception formulation treats sensing as a controlled process~\cite{bajcsy1988}. Motion is chosen partly for the information it produces. In this thesis, the object is not rotated merely to demonstrate dexterity. Rotation changes the direction of gravity in body coordinates and removes a null direction from the CoM regression.

Wilson, Schultz and Murphey study trajectory synthesis for Fisher-information maximization~\cite{wilson2014}. Their work makes the quality of the parameter estimate depend explicitly on the trajectory. The accompanying Baxter study demonstrates the benefit of active parameter identification before executing a task~\cite{wilson2016}. Those studies include dynamic systems. This thesis uses the same information-design principle in a restricted quasi-static setting: the relevant regressor is the gravity-moment balance, and the admissible movement must remain slow and supported by the hands alone.

The resulting information matrix has a simple physical interpretation. Repeating a vertical hang improves precision in two directions but cannot reveal the vertical CoM component. A second gravity direction changes the observable subspace. The decisive quantity is the diversity of accepted measurements, not the commanded tilt amplitude or the number of samples in the file.

\section{Learning coordination in task frames}

Calinon's tutorial develops task-parameterized movement representations in which observations are encoded in several reference frames and adapted to a new task~\cite{calinon2016}. A Gaussian that is narrow in one direction contributes strong precision in that direction after transformation. Products of Gaussians combine these directional constraints into a common representation.

The thesis uses this mechanism for wrenches rather than for a prescribed motion trajectory. Gravity direction is the conditioning input; the two wrist wrenches are the output. The conditional mean expresses a typical allocation, and the conditional covariance expresses how much variation the demonstrations contain around it. The optimizer can deviate from the mean when equilibrium and friction require it, with a penalty determined by the precision.

This use of covariance requires care. Small observed variance can indicate coordinated behaviour, but it can also arise from an apparatus that suppresses a degree of freedom. Valero-Cuevas et al. emphasize the relationship between mechanics, task constraints and interpretations of neuromuscular models~\cite{valero2009}. That relationship motivates the recording checks in this thesis: a torque covariance must be interpreted with the stiffness and transmission of the interface that produced it.

\section{Convex force allocation}

Boyd and Vandenberghe provide the convex-optimization background needed here~\cite{boyd2004}. The contact-force problem is convex when positions, normals, friction coefficients and the desired wrench are fixed during a solve. Equilibrium is affine in contact forces. Each Coulomb cone is a second-order cone. A positive-semidefinite quadratic cost can be represented through a conic epigraph.

Convexity is a property of this local problem. Choosing a contact location, moving a pad to another facet or changing the object pose is not included as a simultaneous convex decision. Contact search and trajectory execution occur outside the allocator and provide a new local problem at the next update.

This distinction prevents an overstatement of robustness. The solver can find an optimum within a given contact model. It cannot ensure that those contacts will persist after the wrist deflects. Nor can it establish that a pose controller will realize eighteen independently allocated force components. Those questions require mechanical and execution measurements.

\section{Compliance and assembly}

Whitney's analysis of quasi-static assembly explains how geometric error, friction and compliance interact during insertion~\cite{whitney1982}. A compliance centre near the contact region can make an insertion error generate a useful corrective movement. A poorly located or poorly oriented compliance can instead increase wedging.

The passive tri-pad follows the same design question at a different scale: where should deformation occur, and which directions should remain stiff? The model separates axial support, surface-conforming tilt and torsion about the grasp axis. These are independent physical channels even when they share one tool.

Herder's work on zero-free-length springs provides a mechanical background for spring arrangements whose deformation energy is written directly as a quadratic displacement~\cite{herder2001}. Guest's treatment of prestressed frameworks shows why stiffness must be derived from the mechanics rather than inferred from a spring constant alone~\cite{guest2006}. Here the small-displacement stiffness is computed from the spring attachment Jacobians. The approximation is stated explicitly; zero free length does not make finite rigid-body rotations linear.

\section{Contact-rich trajectory learning}

Yang et al. formulate crowded-shelf insertion on a quasi-static equilibrium manifold and use an intrinsic haptic metric to evaluate interaction~\cite{yang2025haptic}. Their contact proxies and smooth object representation permit derivatives of the potential. Dynamic movement primitives parameterize the command trajectory, while black-box search identifies a useful insertion route.

The related multi-contact planning work includes stability and noise robustness in DMP-based manipulation~\cite{yang2025multicontact}. Its relevance is the need to evaluate a route under contact and sensing variation. In the present application, the insertion method is retained, while the environment is changed from uniform book neighbours to tetromino profiles and the execution is evaluated with the actual finite tri-pad grasp.

The key comparison is between environments used for training. A trajectory tested with an ideal force applied directly to the object may rely on support unavailable to the hands. A trajectory tested with a shifted starting point may still assume a perfectly known grasp during execution. Both mismatches are examined in the final application.

\section{Position of the present work}

The thesis does not seek a new statistical model, conic solver or DMP formulation. It studies the interfaces between them. Haptic recording must preserve the signal; estimation must identify only observable parameters; optimization must penalize internal forces and respect finite actuation; and training must expose the trajectory to the same grasp errors and contact mechanics as execution. This provides a coherent interpretation of the measured failures as well as the successful runs.


